% ============================================================
%  Part II — Articles, possessives and case
% ============================================================
\gpart{II}{Articles \& Case}

% ------------------------------------------------------------
\gsec{The four cases at a glance}

\begin{gtbl}{colspec={X[0.85,l] X[1.35,l] X[0.75,l] X[1.5,l]}}
  Case & Job in the sentence & Asks & Example \\
  Nominativ & subject & wer? was? & \textbf{Der Hund} beißt den Mann. \\
  Akkusativ & direct object & wen? was? & Der Hund beißt \textbf{den Mann}. \\
  Dativ & indirect object & wem? & Ich gebe \textbf{dem Mann} das Buch. \\
  Genitiv & possessor & wessen? & das Auto \textbf{des Mannes} \\
\end{gtbl}

\tcap{What actually triggers a case}

\begin{flattbl}{colspec={X[1,l] X[2.1,l]}}
  Nominativ & the subject; and after \textit{sein, werden, bleiben, heißen} \\
  Akkusativ & most direct objects; accusative prepositions; time spans \\
  Dativ     & indirect objects; dative verbs; dative prepositions \\
  Genitiv   & possession; genitive prepositions \\
\end{flattbl}

\begin{gnote}
Notice in the examples above that the two nouns swap case, not position:
\textit{Den Mann beißt der Hund} still means the dog does the biting. German
marks roles with endings, so word order is freer than in English.
\end{gnote}

\tcap{The shape of the whole system}

\begin{gendertbl}{}
  Definite & maskulin & feminin & neutrum & Plural \\
  Nominativ & der & die & das & die \\
  Akkusativ & \en{den} & die & das & die \\
  Dativ     & \en{dem} & \en{der} & \en{dem} & \en{den}\,+n \\
  Genitiv   & \en{des}\,+s & \en{der} & \en{des}\,+s & \en{der} \\
\end{gendertbl}

\begin{gnote}
Two housekeeping rules that are easy to forget. \textbf{Dative plural:} the noun
itself takes an extra \textit{-n} (\textit{den Kinder\textbf{n}}) unless the
plural already ends in \textit{-n} or \textit{-s}. \textbf{Genitive masc./neut.:}
the noun takes \textit{-s} or \textit{-es} (\textit{des Hund\textbf{es}}).
\end{gnote}

% ------------------------------------------------------------
\gsec{Nominative}

\begin{gendertbl}{}
  Nominativ & maskulin & feminin & neutrum & Plural \\
  the \eng{definite}   & der    & die    & das    & die    \\
  a \eng{indefinite}   & ein    & eine   & ein    & \xx    \\
  not a \eng{negative} & kein   & keine  & kein   & keine  \\
  my \eng{1st sg}      & mein   & meine  & mein   & meine  \\
  your \eng{2nd sg}    & dein   & deine  & dein   & deine  \\
  your \eng{formal}    & Ihr    & Ihre   & Ihr    & Ihre   \\
  his \eng{3rd sg m}   & sein   & seine  & sein   & seine  \\
  her \eng{3rd sg f}   & ihr    & ihre   & ihr    & ihre   \\
  its \eng{3rd sg n}   & sein   & seine  & sein   & seine  \\
  our \eng{1st pl}     & unser  & unsere & unser  & unsere \\
  your \eng{2nd pl}    & euer   & eure   & euer   & eure   \\
  their \eng{3rd pl}   & ihr    & ihre   & ihr    & ihre   \\
\end{gendertbl}

\begin{gnote}
Every row below the first two carries the \textbf{same endings as \textit{ein}}:
nothing, \textit{-e}, nothing, \textit{-e}. Learn one pattern and the other ten
rows come free. Note that \textit{euer} drops its second \textit{e} whenever an ending
follows: \textit{eure}, not \textit{euere}.
\end{gnote}

\tcap{Examples}

\begin{extbl}{colspec={X[1,l] X[1.6,l]}}
  \textbf{Der} Hund schläft. & The dog is sleeping. \\
  \textbf{Meine} Schwester heißt Anna. & My sister is called Anna. \\
  Das ist \textbf{ein} gutes Buch. & That is a good book. \\
  \textbf{Kein} Mensch weiß das. & Nobody knows that. \\
\end{extbl}

\begin{gnote}
After \textit{sein}, \textit{werden}, \textit{bleiben} and \textit{heißen} the
noun that follows stays nominative — it is not an object:
\textit{Er ist \textbf{ein guter Lehrer}.}
\end{gnote}

% ------------------------------------------------------------
\gsec{Accusative}

\begin{gendertbl}{}
  Akkusativ & maskulin & feminin & neutrum & Plural \\
  the \eng{definite}   & \en{den}     & die    & das    & die    \\
  a \eng{indefinite}   & \en{einen}   & eine   & ein    & \xx    \\
  not a \eng{negative} & \en{keinen}  & keine  & kein   & keine  \\
  my \eng{1st sg}      & \en{meinen}  & meine  & mein   & meine  \\
  your \eng{2nd sg}    & \en{deinen}  & deine  & dein   & deine  \\
  your \eng{formal}    & \en{Ihren}   & Ihre   & Ihr    & Ihre   \\
  his \eng{3rd sg m}   & \en{seinen}  & seine  & sein   & seine  \\
  her \eng{3rd sg f}   & \en{ihren}   & ihre   & ihr    & ihre   \\
  its \eng{3rd sg n}   & \en{seinen}  & seine  & sein   & seine  \\
  our \eng{1st pl}     & \en{unseren} & unsere & unser  & unsere \\
  your \eng{2nd pl}    & \en{euren}   & eure   & euer   & eure   \\
  their \eng{3rd pl}   & \en{ihren}   & ihre   & ihr    & ihre   \\
\end{gendertbl}

\begin{gnote}
\textbf{Only the masculine column changes.} Feminine, neuter and plural are
letter-for-letter identical to the nominative page. If a sentence has no
masculine singular noun in it, the accusative costs you nothing.
\end{gnote}

\tcap{Examples}

\begin{extbl}{colspec={X[1,l] X[1.6,l]}}
  Ich sehe \textbf{den} Hund. & I see the dog. \\
  Sie kauft \textbf{einen} Tisch. & She is buying a table. \\
  Wir haben \textbf{keinen} Hunger. & We are not hungry. \\
  Er liest \textbf{seine} Zeitung. & He is reading his newspaper. \\
\end{extbl}

\begin{gnote}
The accusative also marks duration and definite time without a preposition:
\textit{Ich bleibe \textbf{einen Monat}} — \textit{Er kommt \textbf{jeden Tag}}
— \textit{\textbf{Nächste Woche} fahre ich weg}.
\end{gnote}

% ------------------------------------------------------------
\gsec{Dative}

\begin{gendertbl}{}
  Dativ & maskulin & feminin & neutrum & Plural \\
  the \eng{definite}   & \en{dem}     & \en{der}     & \en{dem}     & \en{den}     \\
  a \eng{indefinite}   & \en{einem}   & \en{einer}   & \en{einem}   & \xx          \\
  not a \eng{negative} & \en{keinem}  & \en{keiner}  & \en{keinem}  & \en{keinen}  \\
  my \eng{1st sg}      & \en{meinem}  & \en{meiner}  & \en{meinem}  & \en{meinen}  \\
  your \eng{2nd sg}    & \en{deinem}  & \en{deiner}  & \en{deinem}  & \en{deinen}  \\
  your \eng{formal}    & \en{Ihrem}   & \en{Ihrer}   & \en{Ihrem}   & \en{Ihren}   \\
  his \eng{3rd sg m}   & \en{seinem}  & \en{seiner}  & \en{seinem}  & \en{seinen}  \\
  her \eng{3rd sg f}   & \en{ihrem}   & \en{ihrer}   & \en{ihrem}   & \en{ihren}   \\
  its \eng{3rd sg n}   & \en{seinem}  & \en{seiner}  & \en{seinem}  & \en{seinen}  \\
  our \eng{1st pl}     & \en{unserem} & \en{unserer} & \en{unserem} & \en{unseren} \\
  your \eng{2nd pl}    & \en{eurem}   & \en{eurer}   & \en{eurem}   & \en{euren}   \\
  their \eng{3rd pl}   & \en{ihrem}   & \en{ihrer}   & \en{ihrem}   & \en{ihren}   \\
\end{gendertbl}

\begin{gnote}
The signature is \textit{-m, -r, -m, -n}. Masculine and neuter are always
identical in the dative, in every single row.
\end{gnote}

\begin{gnote}
\textbf{Dative plural takes an extra -n on the noun itself:}
\textit{die Kinder} $\rightarrow$ \textit{mit den Kinder\textbf{n}};
\textit{die Bücher} $\rightarrow$ \textit{in den Bücher\textbf{n}}. Skip it if
the plural already ends in \textit{-n} or \textit{-s} (\textit{den Autos}).
\end{gnote}

\tcap{Examples}

\begin{extbl}{colspec={X[1,l] X[1.4,l]}}
  Ich gebe \textbf{dem} Kind einen Ball. & I give the child a ball. \\
  Sie hilft \textbf{ihrer} Mutter. & She helps her mother. \\
  Wir fahren mit \textbf{dem} Zug. & We are going by train. \\
  Das gehört \textbf{meinen} Eltern. & That belongs to my parents. \\
\end{extbl}

% ------------------------------------------------------------
\gsec{Genitive}

\begin{gendertbl}{}
  Genitiv & maskulin & feminin & neutrum & Plural \\
  the \eng{definite}   & \en{des}     & \en{der}     & \en{des}     & \en{der}     \\
  a \eng{indefinite}   & \en{eines}   & \en{einer}   & \en{eines}   & \xx          \\
  not a \eng{negative} & \en{keines}  & \en{keiner}  & \en{keines}  & \en{keiner}  \\
  my \eng{1st sg}      & \en{meines}  & \en{meiner}  & \en{meines}  & \en{meiner}  \\
  your \eng{2nd sg}    & \en{deines}  & \en{deiner}  & \en{deines}  & \en{deiner}  \\
  your \eng{formal}    & \en{Ihres}   & \en{Ihrer}   & \en{Ihres}   & \en{Ihrer}   \\
  his \eng{3rd sg m}   & \en{seines}  & \en{seiner}  & \en{seines}  & \en{seiner}  \\
  her \eng{3rd sg f}   & \en{ihres}   & \en{ihrer}   & \en{ihres}   & \en{ihrer}   \\
  its \eng{3rd sg n}   & \en{seines}  & \en{seiner}  & \en{seines}  & \en{seiner}  \\
  our \eng{1st pl}     & \en{unseres} & \en{unserer} & \en{unseres} & \en{unserer} \\
  your \eng{2nd pl}    & \en{eures}   & \en{eurer}   & \en{eures}   & \en{eurer}   \\
  their \eng{3rd pl}   & \en{ihres}   & \en{ihrer}   & \en{ihres}   & \en{ihrer}   \\
\end{gendertbl}

\begin{gnote}
The signature is \textit{-s, -r, -s, -r}. Masculine and neuter match; feminine
and plural match. \textbf{The masculine and neuter noun takes an ending too:}
\textit{-es} after one syllable or after \textit{s, ß, z, tz, x}
(\textit{des Kind\textbf{es}}, \textit{des Hau\textbf{ses}}), otherwise
\textit{-s} (\textit{des Lehrer\textbf{s}}).
\end{gnote}

\tcap{Examples}

\begin{extbl}{colspec={X[1,l] X[1.3,l]}}
  das Auto \textbf{des} Mannes & the man's car \\
  die Farbe \textbf{der} Tür & the colour of the door \\
  wegen \textbf{des} schlechten Wetters & because of the bad weather \\
  trotz \textbf{ihrer} Probleme & despite her problems \\
\end{extbl}

\begin{gnote}
In everyday speech the genitive is usually replaced by \textit{von + Dativ}:
\textit{das Auto \textbf{von dem} Mann} / \textit{vom Mann}. Keep the genitive
for writing, fixed phrases and after genitive prepositions.
\end{gnote}

% ------------------------------------------------------------
\gsec{der-words and ein-words}

Every determiner belongs to one of two families. Which family a word is in tells
you both its own ending and the adjective ending that follows it.

\tcap{der-words \eng{take the definite-article endings}}

\begin{flattbl}{colspec={X[0.85,l] X[1,l] X[0.85,l] X[1.1,l]}}
  dieser  & this            & jeder    & each, every \eng{sg} \\
  jener   & that \eng{lit.} & mancher  & some, many a \\
  welcher & which           & solcher  & such \\
  aller   & all             & derselbe & the same \\
\end{flattbl}

\begin{gendertbl}{}
  dieser & maskulin & feminin & neutrum & Plural \\
  Nominativ & dieser & diese  & dieses & diese  \\
  Akkusativ & diesen & diese  & dieses & diese  \\
  Dativ     & diesem & dieser & diesem & diesen \\
  Genitiv   & dieses & dieser & dieses & dieser \\
\end{gendertbl}

\tcap{ein-words \eng{three slots have no ending at all}}

\begin{flattbl}{colspec={X[0.85,l] X[2.15,l]}}
  Members & ein, kein, and every possessive: mein, dein, sein, ihr, unser, euer, Ihr \\
\end{flattbl}

\begin{gnote}
The two families differ in exactly three cells: masculine nominative, and neuter
nominative and accusative. There a der-word shows \textit{-er} / \textit{-es}
while an ein-word shows nothing (\textit{dies\textbf{er} Mann} vs
\textit{ein Mann}). Everywhere else the endings are identical.
\end{gnote}
